% Ресепшний хүснэгтийн өгөгдлийн бүтэц — UML классын диаграмм (TikZ).
% Эх: reception_tables.py-ийн Room, Service, Staff, ClassGroup, Lesson класс (2026.11.12;
% Staff.email-ийг ресепшний хэрэгцээгээр нэмсэн). Талбарын нэрийг Монголоор бичив.
% 2026.11.23 (багш: «граф дээрх тоонууд хэрэг байна уу, доор нь тайлбар бий»): холбоосын олонлогийн тоог
% (1, 0..1, 0..*) хасав; холбоо бүрийн утгыг шошго ба 3.3-ын бичвэр илэрхийлнэ.
\begin{tikzpicture}[
    font=\footnotesize,
    cls/.style={draw=nmGreenD, line width=0.7pt, rectangle split, rectangle split parts=2, align=left,
                text width=3.6cm, inner sep=3pt, rectangle split part fill={nmGreenS, nmGreen!35}},
    rel/.style={semithick},
    lab/.style={font=\scriptsize, inner sep=1.5pt}]

\node[cls] (room) at (0,0) {\centering\textbf{Өрөө}
  \nodepart{second} дугаар\\давхар\\зориулалт\\тэнхим};
\node[cls] (service) at (7.2,0) {\centering\textbf{Алба, үйлчилгээ}
  \nodepart{second} нэр\\ажлын цаг\\албаны холбоо барих};
\node[cls] (staff) at (-7.2,0) {\centering\textbf{Багш, ажилтан}
  \nodepart{second} нэр\\нэгж, тэнхим\\албан тушаал\\уулзах цаг\\сургуулийн имэйл};
\node[cls] (group) at (-7.2,-5.2) {\centering\textbf{Анги бүлэг}
  \nodepart{second} бүлгийн нэр\\курс\\тэнхим};
\node[cls] (lesson) at (0,-5.2) {\centering\textbf{Хичээл (хуваарийн мөр)}
  \nodepart{second} гараг\\эхлэх, дуусах цаг\\хичээлийн нэр};

% --- холбоос
\draw[rel] (room.east) -- node[lab, below] {байрлуулна} (service.west);

\draw[rel] (staff.east) -- node[lab, below] {ажлын өрөө} (room.west);

\draw[rel] (group.north) -- node[lab, right] {ангийн багш} (staff.south);

\draw[rel] (group.east) -- node[lab, below] {суралцана} (lesson.west);

\draw[rel] (lesson.north) -- node[lab, right] {орно} (room.south);

\draw[rel] (lesson.north west) -- node[lab, sloped, above] {заана} (staff.south east);

\draw[rel] (group.south) -- ++(0,-0.6) -| node[lab, pos=0.25, below] {үндсэн өрөө} ([xshift=3.4cm]lesson.south) |- ([yshift=-3mm]room.east);
\end{tikzpicture}
